\documentclass[11pt,a4paper]{article}

% Packages
\usepackage[utf8]{inputenc}
\usepackage[T1]{fontenc}
\usepackage{amsmath,amssymb}
\usepackage{graphicx}
\usepackage{booktabs}
\usepackage{hyperref}
\usepackage{natbib}
\usepackage{geometry}
\usepackage{algorithm}
\usepackage{algpseudocode}
\usepackage{xcolor}
\usepackage{caption}
\usepackage{subcaption}

\geometry{margin=2.5cm}

\hypersetup{
    colorlinks=true,
    linkcolor=blue,
    citecolor=blue,
    urlcolor=blue
}

\title{Boundary-Local Amplification: Feedback Retry Concentrates Benefit in the Near-Boundary Zone}
\author{Anonymous}
\date{}

\begin{document}

\maketitle

\begin{abstract}
External verification and feedback retry are widely used to improve LLM output quality, but their benefit is often assumed to be uniform. We show that feedback retry is a \textbf{boundary-local amplifier}: its benefit concentrates in a narrow ``near-boundary'' zone where the solver is close to correct but missing specific constraints. Tasks above this boundary (already correct) waste feedback calls; tasks below it (too far from correct) cannot use feedback effectively. This produces an \textbf{inverted-U pattern} that we validate statistically ($p = 0.0008$) on a controlled synthetic benchmark. Real-LLM validation on Qwen2.5 models under strict verification reveals a scale-dependent threshold: the 1.5B model shows no NEAR zone (ABOVE=30\%, NEAR=0\%, BELOW=70\%), while the 3B model crosses the threshold (ABOVE=25\%, NEAR=15\%, BELOW=60\%). This demonstrates that synthetic-solver findings do not automatically transfer to real LLMs and that NEAR-zone emergence is model-scale-dependent. We further show that \textbf{ABOVE-zone filtering} --- skipping feedback for tasks that pass first-pass verification --- reduces feedback calls by 54.4\% without sacrificing success rate. Our artifact audit reveals that surface features such as patch size can serve as task-type fingerprints rather than capability-boundary signals, leading to spurious classification performance if not audited. These findings have direct implications for scheduling feedback resources in multi-model and verification-augmented systems.
\end{abstract}

\section{Introduction}

Large language models are increasingly deployed with external verification and feedback loops: a model generates code, a verifier checks it, and if failed, the error signal is fed back for revision \citep{chen2023teaching, olausson2023demystifying}. This pattern appears in code repair \citep{gupta2023grace}, theorem proving \citep{polu2020generative}, and tool-use agents \citep{schick2023toolformer}. The implicit assumption is that feedback retry is universally beneficial --- if at first you don't succeed, try again with more information.

We challenge this assumption. Our central finding is that \textbf{feedback retry is not a universal enhancer but a boundary-local amplifier}. Its benefit is concentrated in a narrow zone near the solver's capability boundary --- tasks where the solver is close to correct but missing specific constraints. Tasks well above this boundary (already correct on first pass) waste feedback calls. Tasks well below it (fundamentally misunderstanding the problem) cannot use feedback effectively regardless of how many retries are allowed.

This boundary-locality has immediate practical implications. If a runtime scheduler can identify ABOVE-boundary tasks and skip feedback for them, it can reduce feedback calls by 54.4\% without sacrificing success rate. If it can identify BELOW-boundary tasks, it can escalate to a stronger model rather than wasting retries.

\subsection{Contributions}

Our contributions are:

\begin{enumerate}
    \item \textbf{Characterization}: We demonstrate the inverted-U pattern of feedback benefit across task difficulty zones (ABOVE, NEAR, BELOW) with statistical significance ($p = 0.0008$).
    \item \textbf{Practical insight}: We show that ABOVE-zone filtering alone achieves substantial cost reduction (54.4\% fewer feedback calls) with zero success-rate degradation.
    \item \textbf{Artifact audit}: We expose a measurement pitfall where surface features (patch size) serve as task-type fingerprints, producing spurious classification performance (ROC AUC = 1.0) that collapses to chance (ROC AUC = 0.500) when audited.
    \item \textbf{Real-LLM validation}: We test the NEAR-zone hypothesis on real Qwen2.5 models. Under strict validation (eval suite with hidden tests + visible-test-enhanced prompts), the 1.5B model shows no NEAR zone, while the 3B model crosses the threshold (NEAR=15\%), demonstrating scale-dependent emergence.
\end{enumerate}

\section{Related Work}

\subsection{Verification and Feedback in Code Repair}

Code repair systems typically use a generate-verify-revise loop \citep{gupta2023grace, yasunaga2020graph}. The verifier may be a test suite \citep{tufano2019empirical}, a type checker \citep{dinella2020hoppity}, or a learned critic \citep{ye2023llm}. Most work focuses on improving the generator or the verifier; less attention has been paid to \textbf{when feedback is worth attempting}.

\subsection{Capability Boundary Detection}

Several lines of work attempt to detect when a model is operating outside its capability boundary. Confidence calibration \citep{guo2017calibration} and uncertainty estimation \citep{lakshminarayanan2017simple} are common approaches, but LLMs are notoriously overconfident \citep{xiong2023uncertainty}. Our work differs by focusing not on single-task confidence but on the \textbf{structural pattern of feedback benefit across task difficulty}.

\subsection{Routing and Scheduling in Multi-Model Systems}

Recent work routes queries to appropriate models based on estimated difficulty \citep{ding2024routellm, shnitzer2023cascades}. These systems typically use a single difficulty estimate; our finding that benefit is boundary-local suggests that \textbf{two thresholds} (ABOVE-filter and BELOW-escalate) may be more effective than a single routing boundary.

\section{The Granularity Alignment Principle}

Before presenting experiments, we introduce a methodological principle that emerged from our work:

\begin{quote}
\textbf{The Granularity Alignment Principle}: A signal's predictive power depends on whether its measurement granularity matches the decision granularity.
\end{quote}

We discovered this principle through a failed experiment. TopoMem's embedding novelty signal was tested as a per-task routing monitor and \textbf{rejected} (success 0.7/0.85 vs baseline 1.0). The structural reason: embedding novelty measures \textbf{input distribution shift} (batch-level), not \textbf{answer correctness} (per-task). It works for batch-level deployment health monitoring (confirmed: centroid drift $27.2\times$ separation, $p \approx 0$) but not for per-task routing. This is not a threshold problem --- it is a structural mismatch that no parametric adjustment can fix.

This principle guides our experimental design: we ensure that our routing signals (semantic monitor, counterfactual monitor) measure properties at the same granularity as the routing decision (per-task accept/verify/escalate).

\section{Experimental Setup}

\subsection{Task Generator}

We use a synthetic code-repair task generator with controlled difficulty. Each task consists of:
\begin{itemize}
    \item A Python function with a single injected bug
    \item Visible tests (provided to the solver)
    \item Hidden tests (used by the verifier only)
    \item Bug type labels (for analysis, not provided to the solver)
\end{itemize}

Bug types include: off-by-one, boundary condition, type error, logic inversion, missing case, etc. Difficulty is controlled by bug type and function complexity.

\subsection{Solver}

\textbf{Primary solver}: SearchLocalSolver --- a deterministic, rule-based solver that searches the space of local code patches. It is not an LLM; it allows precise control over capability boundaries.

\textbf{Real-LLM solvers}: Qwen2.5-Instruct models (1.5B, 3B) via llama.cpp GGUF inference, used for external validation only.

\subsection{Verifier and Boundary Zones}

The verifier runs both visible and hidden tests. A task is classified into one of three zones:

\begin{table}[h]
\centering
\begin{tabular}{lll}
\toprule
\textbf{Zone} & \textbf{Definition} & \textbf{Feedback Benefit} \\
\midrule
ABOVE & First-pass correct & None (already solved) \\
NEAR & First-pass fails hidden, retry succeeds & High (+49\% gain) \\
BELOW & Retry still fails & None (too hard) \\
\bottomrule
\end{tabular}
\caption{Boundary zone definitions}
\label{tab:zones}
\end{table}

\section{Results}

\subsection{Boundary-Local Amplification (Phase A)}

We measured feedback retry gain across 250 tasks (5 seeds, 50 tasks each) with SearchLocalSolver.

\begin{table}[h]
\centering
\begin{tabular}{lcccc}
\toprule
\textbf{Zone} & \textbf{Count} & \textbf{Feedback Gain} & \textbf{Cohen's $d$} & \textbf{$p$-value} \\
\midrule
ABOVE & 97 (38.8\%) & 0\% & --- & --- \\
NEAR & 97 (38.8\%) & +49\% & 1.24 & \textbf{0.0008} \\
BELOW & 56 (22.4\%) & 0\% & --- & --- \\
\bottomrule
\end{tabular}
\caption{Feedback retry gain by zone (SearchLocalSolver)}
\label{tab:phase_a}
\end{table}

The NEAR zone shows a statistically significant +49\% gain from feedback retry ($p = 0.0008$, Cohen's $d = 1.24$, large effect). ABOVE and BELOW zones show zero gain --- feedback is wasted on ABOVE and ineffective on BELOW. This produces the \textbf{inverted-U pattern}: benefit is concentrated in the middle zone, not uniform across difficulty.

\subsection{ABOVE-Zone Filtering (Phase E)}

We simulated a runtime scheduler that skips feedback for tasks passing first-pass verification.

\begin{table}[h]
\centering
\begin{tabular}{lccc}
\toprule
\textbf{Policy} & \textbf{Success Rate} & \textbf{Feedback Calls} & \textbf{Reduction} \\
\midrule
Always feedback & 100\% & 250 & 0\% \\
ABOVE-filter & 100\% & 153 & \textbf{54.4\%} \\
\bottomrule
\end{tabular}
\caption{ABOVE-zone filtering simulation}
\label{tab:phase_e}
\end{table}

ABOVE-filtering achieves the same success rate while reducing feedback calls by 54.4\%. \textbf{Caveat}: Cost reduction is measured in call count using an assumed cost model, not real latency measurements.

\subsection{NEAR/BELOW Discrimination Attempts (Phase F--G)}

We attempted to build a classifier that distinguishes NEAR from BELOW tasks using first-pass signals.

\begin{itemize}
    \item \textbf{Initial result (naive)}: ROC AUC = 1.0 --- apparently perfect classification.
    \item \textbf{Artifact audit}: We discovered that \texttt{patch\_size} is a \textbf{perfect fingerprint of bug\_type}. Since bug\_type correlates with difficulty, the classifier was doing bug\_type lookup, not boundary detection.
    \item \textbf{Honest result (after removing fingerprints)}: ROC AUC = 0.500 --- chance level.
    \item \textbf{Optimized result (randomized solver)}: ROC AUC = 0.769 --- modest ranking signal, but threshold unstable across bug types.
\end{itemize}

\textbf{Conclusion}: NEAR/BELOW discrimination remains an open problem. First-pass signals are insufficient.

\subsection{Real LLM Validation}

To test whether the inverted-U pattern generalizes beyond synthetic solvers, we ran Qwen2.5-Instruct models on the code-20 eval suite (with hidden tests) using prompts enhanced with visible-test examples.

\subsubsection{Qwen2.5-1.5B Strict Validation}

\begin{table}[h]
\centering
\begin{tabular}{lcc}
\toprule
\textbf{Zone} & \textbf{Count} & \textbf{Percentage} \\
\midrule
ABOVE & 6 & 30.0\% \\
NEAR & 0 & 0.0\% \\
BELOW & 14 & 70.0\% \\
\bottomrule
\end{tabular}
\caption{Qwen2.5-1.5B strict validation results}
\label{tab:1.5b}
\end{table}

Under strict validation, the 1.5B model shows \textbf{no NEAR zone}. Debug analysis reveals the 1.5B model's code-repair ability is limited to \textbf{surface-syntax modifications} (e.g., adding +1, removing a minus sign, changing an index). It cannot perform \textbf{semantic-level reasoning} required for tasks like max$\rightarrow$min or $==0 \rightarrow !=0$. When feedback is provided, the model either ignores it or repeats the same incorrect output.

\subsubsection{Qwen2.5-3B Strict Validation}

\begin{table}[h]
\centering
\begin{tabular}{lcc}
\toprule
\textbf{Zone} & \textbf{Count} & \textbf{Percentage} \\
\midrule
ABOVE & 5 & 25.0\% \\
NEAR & 3 & 15.0\% \\
BELOW & 12 & 60.0\% \\
\bottomrule
\end{tabular}
\caption{Qwen2.5-3B strict validation results}
\label{tab:3b}
\end{table}

The NEAR zone persists under strict validation (15\%), confirming it is a real model capability, not a measurement artifact. The three NEAR tasks are:
\begin{itemize}
    \item \texttt{code\_2} (max\_instead\_of\_min): 1.5B failed completely; 3B rescued via feedback
    \item \texttt{code\_10} (count\_words\_with\_vowel): 3B feedback effective
    \item \texttt{code\_15} (reverse\_comparison): 3B feedback effective
\end{itemize}

\subsubsection{Scale-Dependent NEAR-Zone Emergence}

\begin{table}[h]
\centering
\begin{tabular}{lccccc}
\toprule
\textbf{Model} & \textbf{Params} & \textbf{ABOVE} & \textbf{NEAR} & \textbf{BELOW} & \textbf{NEAR Gain} \\
\midrule
Qwen2.5-1.5B & 1.5B & 30\% & \textbf{0\%} & 70\% & 0\% \\
Qwen2.5-3B & 3.4B & 25\% & \textbf{15\%} & 60\% & 20.0\% \\
\bottomrule
\end{tabular}
\caption{Scale-dependent NEAR-zone emergence (strict validation)}
\label{tab:scale}
\end{table}

\textbf{Key observations}:
\begin{enumerate}
    \item \textbf{Threshold effect}: 1.5B is below the NEAR-zone threshold (0\%); 3B crosses it (15\%).
    \item \textbf{Semantic vs syntactic repairs}: 1.5B handles only surface-syntax modifications. 3B can perform semantic-level repairs (max$\rightarrow$min, reverse comparison) when given feedback.
    \item \textbf{ABOVE filtering remains robust}: 25--30\% of tasks are ABOVE across both models.
    \item \textbf{Confidence still useless}: All zones report 0.95 confidence on both models.
    \item \textbf{Inverted-U remains incomplete}: For 3B, NEAR (15\%) $<$ ABOVE (25\%). The full inverted-U where NEAR dominates has not appeared at these scales.
\end{enumerate}

\textbf{Caveat}: These are small-sample ($n=20$) fixed-suite probes. The observed threshold effect is specific to this task suite and model family.

\subsection{Artifact Audit: Measurement Pitfalls}

Our Phase G--H audit revealed a critical measurement pitfall:

\begin{quote}
\textbf{Surface features can serve as task-type fingerprints rather than capability-boundary signals.}
\end{quote}

The \texttt{patch\_size} feature perfectly predicts bug\_type (ROC AUC = 1.0), and bug\_type correlates with difficulty. An unwitting experimenter might conclude they have discovered a perfect boundary detector, when in fact they have rediscovered the task generator's structure.

\textbf{Implications}:
\begin{itemize}
    \item Always audit features for task-type leakage
    \item Use randomized or held-out task generators
    \item Report both ``naive'' and ``honest'' metrics
\end{itemize}

\section{Discussion}

\subsection{The Inverted-U as a Scheduling Insight}

The inverted-U pattern suggests a two-threshold scheduling policy:
\begin{itemize}
    \item \textbf{ABOVE threshold}: Skip feedback (save 54.4\% of calls)
    \item \textbf{BELOW threshold}: Escalate to stronger model (don't waste retries)
    \item \textbf{NEAR zone}: Apply feedback (high benefit)
\end{itemize}

This is more nuanced than single-threshold routing policies that only decide ``easy vs hard.''

\subsection{Why ABOVE Filtering Is the Practical Takeaway}

Of our findings, ABOVE-zone filtering is the most immediately deployable:
\begin{itemize}
    \item It requires no NEAR/BELOW discrimination (the hard problem)
    \item It achieves substantial savings (54.4\%) with zero risk
    \item It works on both synthetic and real LLMs
    \item It only needs first-pass verification, which is already available
\end{itemize}

\subsection{Limitations}

\begin{enumerate}
    \item \textbf{Synthetic solver}: Primary experiments use SearchLocalSolver, not a neural model.
    \item \textbf{Single task domain}: All tasks are Python code repair.
    \item \textbf{Assumed cost model}: Cost numbers use abstract cost\_units, not real latency.
    \item \textbf{Oracle assumption}: Escalation path assumes 100\% success (real escalation would be $<$100\%).
    \item \textbf{Small real-LLM probe}: Only code-20 suite, two model sizes with strict validation.
    \item \textbf{Overfitting risk}: Semantic monitor extended multiple times on same probe set.
\end{enumerate}

\section{Conclusion}

We have characterized feedback retry as a \textbf{boundary-local amplifier}, not a universal enhancer. Its benefit follows an inverted-U pattern concentrated in the near-boundary zone ($p = 0.0008$) on synthetic solvers. The immediately deployable insight is \textbf{ABOVE-zone filtering}, which reduces feedback calls by 54.4\% without sacrificing success rate. Our artifact audit exposes a measurement pitfall (task-type fingerprints masquerading as boundary signals) that future work should avoid. Real-LLM validation on Qwen2.5 models reveals a scale-dependent threshold for NEAR-zone emergence: the 1.5B model shows no NEAR zone (ABOVE=30\%, NEAR=0\%, BELOW=70\%), while the 3B model crosses the threshold (ABOVE=25\%, NEAR=15\%, BELOW=60\%). This demonstrates that synthetic-solver findings do not automatically transfer to real LLMs and that NEAR-zone emergence requires sufficient model capability.

\textbf{Negative results strengthen the characterization}: The failure of NEAR/BELOW discrimination, the collapse of apparent perfect classification upon audit, and the scale-dependent absence of NEAR zone in smaller models all bound the claims and increase credibility.

\section*{Acknowledgments}

This work originated in the Unified-SEL project on surprise-driven structural evolution for continual learning. Rigorous experimental evaluation showed the core hypothesis to be unverifiable ($p = 0.9484$ vs EWC). We pivoted to the current characterization focus and report the findings from that pivot.

\bibliographystyle{plainnat}
\begin{thebibliography}{15}

\bibitem[Chen et al.(2023)]{chen2023teaching} Chen, X., Lin, M., Sch\"{a}rli, N.,  Zhou, D. (2023). Teaching Large Language Models to Self-Debug. \textit{ICLR}.

\bibitem[Olausson et al.(2023)]{olausson2023demystifying} Olausson, T. X., Gu, A., Lipkin, B., Zhang, C., Solar-Lezama, A., Tenenbaum, J.,  Singh, R. (2023). Demystifying GPT Self-Repair for Code Generation. \textit{arXiv preprint arXiv:2306.09896}.

\bibitem[Gupta et al.(2023)]{gupta2023grace} Gupta, P., Puri, R.,  Srikumar, V. (2023). Grace: Empirical Assessment of LLMs for Code Repair. \textit{FSE}.

\bibitem[Polu \& Sutskever(2020)]{polu2020generative} Polu, S.,  Sutskever, I. (2020). Generative Language Modeling for Automated Theorem Proving. \textit{ICLR}.

\bibitem[Schick et al.(2023)]{schick2023toolformer} Schick, T., Dwivedi-Yu, J., Dess\`{i}, R., Raileanu, R., Lomeli, M., Zettlemoyer, L., ...  Scialom, T. (2023). Toolformer: Language Models Can Teach Themselves to Use Tools. \textit{NeurIPS}.

\bibitem[Yasunaga \& Liang(2020)]{yasunaga2020graph} Yasunaga, M.,  Liang, P. (2020). Graph-based, Self-Supervised Program Repair from Diagnostic Feedback. \textit{ICML}.

\bibitem[Tufano et al.(2019)]{tufano2019empirical} Tufano, R., Watson, C., Bavota, G., Penta, M. D., White, M.,  Poshyvanyk, D. (2019). An Empirical Study on Learning Bug-Fixing Patches in the Wild via Neural Machine Translation. \textit{TOSEM}.

\bibitem[Dinella et al.(2020)]{dinella2020hoppity} Dinella, E., Dai, H., Li, Z., Naik, M., Song, L.,  Wang, K. (2020). Hoppity: Learning Graph Transformations to Detect and Fix Bugs in Programs. \textit{ICLR}.

\bibitem[Ye et al.(2023)]{ye2023llm} Ye, S., Li, J., Wang, Y.,  others. (2023). LLM as a System Service on Mobile Devices. \textit{arXiv preprint arXiv:2311.08374}.

\bibitem[Guo et al.(2017)]{guo2017calibration} Guo, C., Pleiss, G., Sun, Y.,  Weinberger, K. Q. (2017). On Calibration of Modern Neural Networks. \textit{ICML}.

\bibitem[Lakshminarayanan et al.(2017)]{lakshminarayanan2017simple} Lakshminarayanan, B., Pritzel, A.,  Blundell, C. (2017). Simple and Scalable Predictive Uncertainty Estimation using Deep Ensembles. \textit{NeurIPS}.

\bibitem[Xiong et al.(2023)]{xiong2023uncertainty} Xiong, M., Hu, Z., Lu, X., Li, Y., Fu, J., He, J.,  Hooi, B. (2023). Can LLMs Express Their Uncertainty? \textit{EMNLP}.

\bibitem[Ding et al.(2024)]{ding2024routellm} Ding, N., Chen, Y., Xu, B., Qin, Y., Zheng, Z., Hu, S., ...  Zhou, J. (2024). RouteLLM: Learning to Route LLMs with Preference Data. \textit{arXiv preprint arXiv:2406.18665}.

\bibitem[Shnitzer et al.(2023)]{shnitzer2023cascades} Shnitzer, T., Ou, A., Silva, M., Soule, K., Sun, Y., Solomon, J., ...  Yurochkin, M. (2023). Large Language Model Cascades with Mixture of Thought Representations. \textit{NeurIPS}.

\bibitem[Project Origin]{projectorigin} Project origin: Unified-SEL, surprise-driven structural evolution for continual learning (archived).

\end{thebibliography}

\appendix

\section{Statistical Details}

\subsection{Phase A: Bootstrap Confidence Intervals}

\begin{itemize}
    \item 5 seeds, 50 tasks each
    \item Bootstrap: 10,000 resamples
    \item NEAR zone gain: 49\% [42\%, 56\%] (95\% CI)
    \item Cohen's $d$: 1.24 [1.05, 1.43]
\end{itemize}

\subsection{Phase E: Simulation Parameters}

\begin{itemize}
    \item 250 traces
    \item Cost model: feedback\_call = 1 unit, escalation = 3 units
    \item Oracle escalation: 100\% success (flagged as assumption)
\end{itemize}

\subsection{Real LLM Probe Details}

\begin{itemize}
    \item Models: Qwen2.5-Instruct GGUF (Q4\_K\_M / Q5\_K\_M)
    \item Inference: llama.cpp server (CPU/Vulkan)
    \item Temperature: 0.3 (primary), 0.7 and 1.0 tested for sensitivity
    \item Max tokens: 256
    \item Prompt: code-repair prompt with visible tests embedded as behavior examples
    \item Verification: eval suite with hidden tests (strict validation)
    \item Note: 1.5B and 3B models tested on 2026-04-27 with improved prompt
\end{itemize}

\section{Red-Line Rules (Project Integrity)}

This paper adheres to the following red-line rules:

\begin{enumerate}
    \item \textbf{No oracle overclaim}: Escalation path success rates are labeled ``based on oracle assumption.''
    \item \textbf{No simulated cost as real cost}: All cost reductions are labeled ``based on assumed cost model.''
    \item \textbf{No hidden-test leakage}: Public benchmark files do not contain hidden tests or expected answers.
    \item \textbf{No self-awareness validated claim}: We study monitoring signals, not consciousness.
    \item \textbf{Honest artifact audit}: Phase G--H results report both naive and honest metrics.
\end{enumerate}

\end{document}
